\documentclass[11pt]{article}

\setlength{\parindent}{0pt}
\usepackage{xltxtra}
\usepackage{hyperref}
\hypersetup{hidelinks}
\usepackage{url}
\urlstyle{tt}
\usepackage{xcolor}
\definecolor{CVBlue}{RGB}{23,110,191}
\usepackage{calc}
\usepackage{graphicx}
\usepackage{tikz}
\usetikzlibrary{calc}
\usepackage{fontspec}
\usepackage{xeCJK}
\usepackage{enumitem}
\CJKsetecglue{}

%% 字体设置
\setmainfont[
    Path = fonts/Main/,
    Extension = .otf,
    BoldFont = texgyretermes-bold.otf,
]{texgyretermes-regular.otf}

\setCJKmainfont[
    Path = fonts/hansans/,
    Extension = .ttf,
    BoldFont = NotoSansSC-Bold.ttf,
]{NotoSansSC-Regular.otf}

\usepackage{relsize}
\usepackage{xspace}
\usepackage{fontawesome}

\usepackage[
    a4paper,
    left=1.2cm,
    right=1.2cm,
    top=1.1cm,      % 进一步压缩上边距
    bottom=0.7cm,   % 进一步压缩下边距
    nohead
]{geometry}

\renewcommand{\baselinestretch}{1.15} % 行间距微调，兼顾紧凑与可读性

\usepackage{titlesec}
\usepackage{enumitem}
\setlist{noitemsep}
\setlist[itemize]{topsep=0.15em, leftmargin=*} % 极致压缩列表间距
\setlist[enumerate]{topsep=0.15em, leftmargin=*}

\newlength{\interProjectSpacing}
\setlength{\interProjectSpacing}{0.3em} % 压缩项目间空白
\newcommand{\projectsep}{\vspace{\interProjectSpacing}}

\newlength{\intraProjectTitleSep}
\setlength{\intraProjectTitleSep}{0.2em} % 压缩标题换行间距
\newcommand{\titlebreak}{\\[\intraProjectTitleSep]}

\newlength{\intraProjectListTopSep}
\setlength{\intraProjectListTopSep}{0.1em} % 压缩列表与上文间距

\titleformat{\section}
{\normalsize\bfseries\raggedright}
{}{0em}
{}
[{\color{CVBlue}\titlerule}]
\titlespacing*{\section}{0cm}{*0.9}{*0.7} % 压缩section标题上下间距

\begin{document}
\pagenumbering{gobble}

\begin{tikzpicture}[remember picture, overlay]
    \node[anchor = north west] at ($(current page.north west)+(0.5cm,+0.7cm)$) {\includegraphics[height=5cm]{zju.png}}; % 微调校徽大小，节省顶部空间
\end{tikzpicture}%

\centerline{\LARGE\bfseries{欧庆辉}}
\centerline{\normalsize{\faPhone\ 152-5117-6009 \quad \faEnvelopeO\ \href{mailto:15251176009@163.com}{15251176009@163.com}}}

\centerline{\normalsize{\textbf{求职意向：}AI应用开发工程师、软件开发工程师}}

\section{\makebox[\widthof{\faGraduationCap}][c]{\color{CVBlue}\faGraduationCap}\ 教育背景}
\begin{tabular}{l l r}
\textbf{浙江大学} & 人工智能/硕士 & 2024.9 -- 2027.6 \\[0.25em] % 压缩表格行间距
\textbf{苏州大学} & 高分子材料与工程/学士 & 2018.9 -- 2022.6 \\
\end{tabular}

% ========== 技术栈章节（完整保留要求内容） ==========
\section{\makebox[\widthof{\faWrench}][c]{\color{CVBlue}\faWrench}\ 技术栈}
\textbf{后端：} Java 8+/21、Spring Boot、MySQL、Redis、MinIO、Kafka、Elasticsearch \\
\textbf{AI/RAG：} RAG、GraphRAG、Ollama、LangGraph4j、Spring AI、Prompt工程

\section{\makebox[\widthof{\faUsers}][c]{\color{CVBlue}\faUsers}\ 项目经历}

% ========== Tiny-RAG 企业知识库系统（精简冗余，保留核心） ==========
\textbf{Tiny-RAG 企业知识库系统} \hfill  \titlebreak
项目描述：设计并实现高可用、高召回的企业级混合检索知识库系统，解决传统RAG在专有名词识别、长文档语义割裂、大文件上传可靠性等痛点，支持多租户安全访问与低Token智能问答。
\begin{itemize}[nosep, topsep=\intraProjectListTopSep]
    \item \textbf{混合检索引擎架构}：基于BM25（稀疏）与句向量Embedding（稠密）构建并行双引擎，结合KNN向量召回与关键词检索，引入Cross-Encoder高精度重排，提升专有名词与长尾语义覆盖率。
    \item \textbf{双引擎索引实现}：基于Elasticsearch+IK分词搭建BM25关键词索引，通义千问2048维Embedding做语义向量化，通过RRF算法融合检索结果并重排。
    \item \textbf{长文档语义优化}：提出Parent-Child Indexing父子索引策略，小粒度子块定位检索、回溯父级恢复上下文，兼顾召回精度与语义完整性。
    \item \textbf{高可靠文件上传架构}：基于Kafka解耦「上传→解析→向量化」全流程，支持分片上传与断点续传；Redis Bitmap跟踪分片状态，MinIO+MD5校验实现可靠合并。
    \item \textbf{智能检索与安全管控}：Dynamic Top-K调度动态调整LLM输入片段数量，降低Token消耗；LLM驱动Query-Rewrite模块；RBAC多级权限控制与租户隔离。
\end{itemize}

\projectsep
\vspace{0.8em}
% ========== PaiAgent 完整保留用户提供的详细文字版 ==========
\textbf{PaiAgent - 企业级AI工作流可视化编排平台} \hfill  \titlebreak
项目描述：基于LangGraph4j与Spring AI的企业级AI工作流平台，支持可视化拖拽编排多模型（OpenAI、DeepSeek、通义千问）与工具节点，DAG引擎按拓扑顺序执行复杂AI任务，实现零代码构建AI Agent应用。
\begin{itemize}[nosep, topsep=\intraProjectListTopSep]
    \item 基于LangGraph4j StateGraph构建工作流引擎，GraphBuilder负责节点注册与边连接，NodeAdapter适配现有执行器为异步节点，StateManager管理状态传递，支持状态图、条件分支的高级编排。
    \item \textbf{设计Spring AI多模型统一接入架构}：ChatClientFactory动态工厂，运行时根据节点配置动态创建OpenAI兼容ChatClient，OpenAiChatOptions统一参数配置，实现多厂商LLM无缝切换。
    \item \textbf{模板方法模式重构LLM节点执行器}：抽象AbstractLLMNodeExecutor基类封装通用流程，子类仅需实现getNodeType()，将5个执行器代码从800+行精简至约10行/个。
    \item \textbf{实现Prompt模板变量替换}：支持input静态值与reference动态引用，从上游节点输出自动获取参数，实现节点间数据灵活映射。
    \item \textbf{实现DAG工作流解析引擎}：基于Kahn算法拓扑排序确定执行顺序，DFS深度优先搜索检测循环依赖防止死锁，支持一对多、多对一节点连接。
    \item \textbf{基于Spring SseEmitter实现实时反馈}：设计ExecutionEvent事件模型，Consumer<ExecutionEvent>回调将LLM流式生成内容实时推送到前端调试面板。
\end{itemize}

\section{\makebox[\widthof{\faTrophy}][c]{\color{CVBlue}\faTrophy}\ 获奖情况}
\begin{itemize}[nosep]
    \item 浙江大学一等学业优秀奖学金 \hfill 2025.11
    \item 唐仲英德育奖学金 \hfill 2019 -- 2022
\end{itemize}

\end{document}
